\documentclass[oneside]{scrbook}

% \usepackage[utf8]{inputenc}
\usepackage[T1]{fontenc}
\usepackage{version}
% \usepackage[latin1]{inputenc}
\usepackage[italian]{babel}
\usepackage{algpseudocode}
\usepackage{amsmath}
\usepackage{booktabs}
\usepackage{color}
\usepackage{amsfonts}
\usepackage{url}
\usepackage{algorithm}
\usepackage{chngcntr}
\usepackage[answerdelayed]{exercise}
\usepackage{graphicx}
\usepackage{caption}
\usepackage{listings}
\usepackage{tikz}
\usepackage{pgf}
\usepackage{colortbl}
\usepackage{listingsutf8}
\usepackage[pdftitle = {Dispense di esercizi}, pdfauthor = {Gianmarco Tedeschi e Alberto Marchesi}, pdfsubject = {Fondamenti di Informatica ($082746$ - A.A.~$2025-2026$)}, pdfcreator = {Gianmarco Tedeschi e Alberto Marchesi}, pdfproducer = {-}]{hyperref}
% \usepackage[automark,nouppercase]{scrpage2}
\usepackage{mdframed}
\usepackage{scrlayer-scrpage} % Necessary package for page styles

%\pagestyle{scrheadings}

%\usepackage{titlesec}
%\newcommand{\sectionbreak}{\clearpage}
\counterwithin{Exercise}{chapter}
\counterwithin{Answer}{chapter}

\algrenewcommand\algorithmicthen{\textbf{Allora}}
\algrenewcommand\algorithmicif{\textbf{Se}}
\algrenewcommand\algorithmicelse{\textbf{Altrimenti}}
\algrenewcommand\algorithmicend{\textbf{Chiudi}}
\algrenewcommand\algorithmicwhile{\textbf{Finch\'e}}
\algrenewcommand\algorithmicdo{\textbf{Esegui}}
\algrenewcommand{\algorithmiccomment}[1]{\hskip3em\texttt{//#1}}

\setlength{\parindent}{0in}
\setlength{\parskip}{10pt}

\newmdenv[%
        outerlinewidth=2,%
        roundcorner=10pt,%
        leftmargin=0,%
        rightmargin=0,%
        outerlinecolor=black,
        innertopmargin=\topskip,%
        splittopskip=\topskip,%
]{algo}


\ohead{\headmark}
\ihead{}
\chead{}
\ifoot{A.A.~$2025-2026$}
\cfoot{Fondamenti di Informatica (082746) - G. Tedeschi}
\ofoot{Pag.~\pagemark}
\setheadsepline{.4pt}
\setfootsepline{.4pt}

\definecolor{lbcolor}{rgb}{0.9,0.9,0.9}

\lstdefinestyle{cstyle}{%
  inputencoding=utf8/latin1,
  breaklines=true,
  breakatwhitespace=false,
  basicstyle=\small\ttfamily,
  language=c,
  frame=single,
  showstringspaces=false
}

\lstdefinestyle{cstylesmall}{%
  inputencoding=utf8/latin1,
  breaklines=true,
  breakatwhitespace=false,
  basicstyle=\scriptsize\ttfamily,
  language=c,
  frame=single,
  showstringspaces=false
}

\lstdefinestyle{console}{%
  inputencoding=utf8/latin1,
  breaklines=true,
  breakatwhitespace=false,
  basicstyle=\scriptsize\ttfamily,
  showstringspaces=false,
  backgroundcolor=\color{lbcolor}
}


\lstnewenvironment{clang}{\lstset{style=cstyle}}{}
\lstnewenvironment{clangsmall}{\lstset{style=cstylesmall}}{}

\addtocounter{chapter}{5}

\begin{document}
\chapter{Struct}
Questa dispensa propone esercizi sulla scrittura di algoritmi, in linguaggio C, utili alla comprensione della definizione di strutture e di tipi, oltre ad un riepilogo sulla codifica dei numeri con complemento a due.

I tipi di dato strutturato sono dichiarabili in C tramite la parola chiave \texttt{struct}
\begin{lstlisting}[style=cstyle]
void main(){

struct {
	tipo_campo1 nome_campo1;
	tipo_campo2 nome_campo2;
	...
} nome_struttura;

istruzioni;

}
\end{lstlisting}		

A questo punto potremo utilizzare \texttt{nome\_struttura} come se fosse una variabile. Per accedere ai campi della struttura si utilizza la sintassi \texttt{nome\_struttura.nome\_campo}. Grazie ad esso possiamo accedere (leggere o scrivere) nei campi della struttura. L'assegnamento tra strutture \`e possibile solo se le strutture sono dello stesso tipo (stessi campi, per nome e tipo, con stesso ordine di dichiarazione). Non \`e possibile fare confronti tra strutture in maniera automatica.
		
Oltre ad i tipi nativi del C (ad esempio \texttt{int}, \texttt{char} o \texttt{float}) \`e possibile dichiarare dei tipi definiti dall'utente tramite la parola chiave \texttt{typedef}. Ad esempio, volendo dichiarare dei vettori della stessa dimensione \texttt{n\_elementi}:

\begin{lstlisting}[style=cstyle]
#define LEN n_elementi

typedef tipo_variabile Nuovo_nome_tipo[LEN];

void main(){

Nuovo_nome_tipo nome_variabile;

istruzioni;
}
\end{lstlisting}
\`E buona norma dichiarare nuovi tipi con nomi che inizino con una \textbf{lettera maiuscola}, con un prefisso o con un suffisso scelto.
		
La tecnica della dichiarazione di tipo si pu\`o combinare con la dichiarazione di strutture. Ad esempio:
\begin{lstlisting}[style=cstyle]
#define LEN n_elementi

typedef tipo_variabile Nome_tipo[LEN][LEN];

typedef struct {
	tipo_campo1 nome_campo1;
	tipo_campo2 nome_campo2;
	...
} Nome_tipo_struttura;

void main(){

Nome_tipo nome_variabile;
Nome_tipo_struttura nome_struttura;

istruzioni;
}
\end{lstlisting}
La definizione di un tipo non implica l'istanziazione di una variabile di quel tipo. La dichiarazione della variabile avverrà di seguito.

\paragraph{Enum}
Gli enumerated types contengono un elenco di costanti che possono essere indirizzate con valori integer. Per dichiarare tali tipi si utilizza \texttt{enum}; vengono dichiarati i tipi e le variabili come nell'esempio che segue:
\begin{lstlisting}[style=cstyle]
	enum colori {rosso, giallo, verde, blu} pennarello;
	enum giorni {lun,mar,mer,gio,ven,sab,dom} settimana;
	enum colori pulsante, nastro;
\end{lstlisting}

In tale esempio viene dichiarato \texttt{colori} come enumerated type e la variabile \texttt{pennarello} con $4$ valori accettabili definiti, mentre la variabile \texttt{settimana} di tipo \texttt{giorni} ha $7$ valori accettabili definiti. Le variabili \texttt{pulsante} e \texttt{nastro} sono di tipo \texttt{colori}. Ogni item nell'elenco di valori accettabili \`e detto enumeration constant. Il C mappa ogni enumeration constant ad un'intero, per cui \`e ad esempio possibile scrivere:
\begin{lstlisting}[style=cstyle]
	settimana = verde;
\end{lstlisting}
che come risultato fa si che settimana abbia valore $2$, perch\`e di default a ogni membro dell'elenco di variabili \`e assegnato un valore incrementale partendo da $0$ per il primo valore. \`E comunque possibile definire valori diversi agli elementi:
\begin{lstlisting}[style=cstyle]
	enum colori {rosso=10, giallo=30, verde, blu=giallo};
\end{lstlisting}

Un ulteriore esempio relativo all'assegnazione di valori diversi \`e il seguente:
\begin{lstlisting}[style=cstyle]
	enum escapes {bell='\a', backspace='\b', tab='\t', newline='\n', vtab='\v', return='\r'};
\end{lstlisting}
\`E anche possibile iniziare ad assegnare i valori da un valore iniziale diverso da $0$:
\begin{lstlisting}[style=cstyle]
	enum months (jan=1,feb,mar,...dec);
\end{lstlisting}
dove \`e implicito che febbraio sia corrispondente al valore $2$, marzo al $3$ e cos\`i via.

				
%\paragraph{Riepilogo: IEEE 754-1985 standard}
%Decimale -> IEEE standard
%\begin{itemize}
		%\item Definisco il bit di segno $S = 0$ per positivo e $S = 1$ per negativo
		%\item Codifico in virgola fissa in base 2, parte frazionaria e parte intera
		%\item Porto il numero binario in forma normalizzata in base 2
		%\item Definisco a 23 bit come la mantissa $M$ senza il primo bit (sempre 1)
		%\item Calcolo l'esponente $E$ (aggiungo 127 a quello che ho trovato durante la normalizzazione)
		%\item Compongo il numero $S/E/M$
%\end{itemize}
%
%IEEE standard -> Decimale
%\begin{itemize}
		%\item Decompongo il numero $S/E/M$
		%\item Calcolo l'esponente $E$ e sottraggo 127
		%\item Sposto la virgola in base all'esponente alla mantissa $M$
		%\item Converto in decimale la parte intera
		%\item Converto in decimale la parte frazionaria
		%\item Definisco il bit di segno $S = 0$ per positivo e $S = 1$ per negativo
%\end{itemize}
	
\section{Esercizi}

\begin{Exercise}[label={ex:capocannoniere}]
		Un campionato di calcio \`e composto da venti squadre. Ciascuna squadra ha un nome (lungo al pi\'u 30 caratteri), un numero di punti accumulati durante la stagione e al più 23 calciatori. Ciascun calciatore ha un cognome (lungo al pi\'u 20 caratteri), un'et\'a, un numero di maglia, un numero di presenze in stagione ed un numero di gol realizzati in stagione.
		\begin{enumerate}
			\item Specificare i tipi Campionato, Squadra e Calciatore.
			\item Viene data una variabile di nome serie\_a di tipo Campionato, gi\'a inizializzata e riempita dall'utente con tutte le squadre del campionato. Scrivere un frammento di codice che stampi il cognome, il numero di gol realizzati e la squadra di appartenenza del calciatore che ha realizzato pi\'u gol nel campionato.
		\end{enumerate}
\end{Exercise}

\begin{Answer}[ref={ex:capocannoniere}]
	\begin{enumerate}
		\item \lstinputlisting[style=cstylesmall]{src/campionato.c}
		
		\item \lstinputlisting[style=cstylesmall]{src/capocannoniere.c}
	\end{enumerate}
\end{Answer}
\begin{Exercise}[label={ex:libretto}]
		Scrivere un programma che permetta all'utente di gestire un libretto di voti di esami, per un massimo di 20 esami. Ogni esame sostenuto ha i seguenti attributi:

		\begin{itemize}
		    \item nome del corso (stringa di massimo 256 caratteri)
		    \item voto (minimo 18, massimo 30)
		    \item data (in formato gg/mm/aaaa)
		    \item codice del corso (codice di massimo 6 cifre)
		\end{itemize}

		Il programma permette all'utente di svolgere le seguenti operazioni:

		\begin{itemize}
		    \item \textbf{inserimento:} inserisce un esame in fondo alla lista degli esami eventualmente presenti.
		    \item \textbf{stampa:} stampa tutti gli esami presenti, con i dettagli di cui sopra.
		    \item \textbf{ricerca:} chiede all'utente di inserire un codice e cerca nel libretto la presenza di un esame corrispondente a quel codice. Se presente, stampa tale esame.
		    \item \textbf{uscita:} esce dal programma.
		\end{itemize}
		tramite un menù di scelta di questo tipo:
		\begin{verbatim}
		    [i] inserimento nuovo esame
		    [s] stampa del libretto
		    [r] ricerca per codice
		    [x] uscita
		\end{verbatim}

		Il programma deve continuare a presentare tale menù, fino a quando l'utente non preme il tasto per uscire.
\end{Exercise}

\begin{Answer}[ref={ex:libretto}]
  	\lstinputlisting[style=cstylesmall]{src/libretto.c}
\end{Answer}
\begin{Exercise}[label={ex:magazzino}]
    Si vuole realizzare un programma per la gestione di un magazzino.
    A tal scopo definiamo un'opportuna struttura \texttt{prodotto} definita da:
    \begin{itemize}
        \item un campo \texttt{id} di tipo intero che rappresenta il codice univoco del prodotto;
        \item un campo \texttt{nome} di tipo stringa (array di caratteri) di lunghezza massima 50 che rappresenta il nome del prodotto;
        \item un campo \texttt{quantita} di tipo float che rappresenta la quantità disponibile del prodotto;
        \item un campo \texttt{prezzo} di tipo float che rappresenta il prezzo del prodotto.
    \end{itemize}

    Si scriva un programma che, dichiarato uno array, di massimo 10 elementi di tipo \texttt{prodotto} (il magazzino), permetta di:
    \begin{itemize}
        \item inserire un nuovo prodotto nel magazzino (chiedendo all'utente di inserire i valori dei campi);
        \item incremetare la quantità di un prodotto esistente (dato il suo \texttt{id});
        \item visualizzare tutti i prodotti presenti nel magazzino con i loro dettagli;
        \item vendere un prodotto (dato il suo \texttt{id} e la quantità da vendere), aggiornando la quantità disponibile;
        \item uscire dal programma.
    \end{itemize}
\end{Exercise}

\begin{Answer}[ref={ex:magazzino}]
  	\lstinputlisting[style=cstylesmall]{src/magazzino.c}
\end{Answer}
\begin{Exercise}[label={ex:azienda}]
		Le seguenti dichiarazioni definiscono tipi di dati relativi alla categoria degli impiegati di un'azienda (gli impiegati possono essere di prima, seconda, ..., quinta categoria), agli uffici occupati da tali impiegati, all'edificio che ospita tali uffici (edificio diviso in 20 piani ognuno con 40 uffici).

		\begin{clang}
/* definizioni dei tipi */
typedef struct {
  char nome[20], cognome[20];
  int cat; // contiene valori tra 1 e 5
  int stipendio;
} Impiegato;

typedef enum{
  nord, nordEst, est, sudEst, sud, sudOvest, ovest, nordOvest
} Esposizione;

typedef struct {
  int superficie; /* in m^2 */
  Esposizione esp;
  Impiegato occupante;
} Ufficio;

/* definizioni delle variabili */
Ufficio torre[20][40];

/* rappresenta un edificio di 20 piani con 40 uffici per piano */
		\end{clang}

\begin{enumerate}
		\item Si scriva un frammento di codice (che includa eventualmente anche le dichiarazioni di ulteriori variabili) che stampi il cognome, lo stipendio e la categoria di tutte e sole le persone che occupano un ufficio orientato a sud oppure a sudEst e avente una superficie compresa tra $20$ e $30$ metri quadri;
  	\item Visualizzi a schermo i piani che non hanno neanche un ufficio esposto a nord;
  	\item Stampi lo stipendio medio degli impiegati in questi piani che si chiamano ``Giacomo'' di nome.
\end{enumerate}
\end{Exercise}

\begin{Answer}[ref={ex:azienda}]
  \lstinputlisting[style=cstylesmall]{src/torre.c}
  \lstinputlisting[style=cstylesmall]{src/torre2.c}
  \lstinputlisting[style=cstylesmall]{src/torre3.c}
\end{Answer}
\begin{Exercise}[label={ex:facebook}]
  \Question Si definiscano le seguenti variabili che specificano le strutture dati per rappresentare i messaggi scambiati attraverso Facebook e ai profili degli utenti:
  		\begin{itemize}
  				\item una struttura che definisca un'utente, con nome, cognome e e-mail, tutti composti da caratteri alfanumerici;
  				\item una struttura che definisca un messaggio, con un contenuto alfanumerico, un mittente scelto tra gli utenti, un numero fissato di destinatari e i destinatari, anch'essi scelti tra gli utenti (fino ad un massimo di 256).
  		\end{itemize}

		\Question Definire, usando il linguaggio C, un'appropriata variabile per memorizzare le informazioni relative a 10 utenti e 10 messaggi. Queste variabili dovranno essere chiamate \texttt{utenti\_facebook} e \texttt{messaggi\_inviati}.

		\Question Scrivere in linguaggio C, aggiungendo eventualmente opportune dichiarazioni di variabili, un frammento di codice che permetta di visualizzare a video l'indirizzo email di tutti gli utenti che hanno spedito almeno un messaggio e che non siano tra i destinatari di tale messaggio. Si presupponga che la variabile \texttt{messaggi\_inviati} sia inizializzata correttamente con le informazioni relative a 10 messaggi. Si presupponga inoltre che l'indirizzo email sia identificatore univoco di un utente: pertanto si può utilizzare l'indirizzo email per verificare se due utenti sono lo stesso utente.
		
\end{Exercise}

\begin{Answer}[ref={ex:facebook}]

\begin{clangsmall}
typedef struct {
  char nome[24];
  char cognome[24];
  char email[64];
} Utente;

typedef struct {
  char contenuto[256];
  Utente mittente;
  int numDestinatari;
  Utente destinatari[256];
} Messaggio;
\end{clangsmall}

\begin{clangsmall}
Utente utenti_facebook[10];
Messaggio messaggi_inviati[10];		
\end{clangsmall}

Gli utenti che hanno spedito almeno un messaggio sono necessariamente presenti come mittenti nella variabile \texttt{messaggi\_inviati}:

\begin{clangsmall}
void main () {
  int i, j;
  Utente utenti_facebook[10];
  Messaggio messaggi_inviati[10];		

  for (i = 0; i < 10; i++) {
    j = 0;

    //ricerca destinatario.email != mittente.email con strcmp()
    while (j < messaggi_inviati[i].numDestinatari &&
              strcmp(messaggi_inviati[i].mittente.email,
                     messaggi_inviati[i].destinatari[j].email) != 0) {
      j++;
    }

    if (j == messaggi_inviati[i].numDestinatari) //non trovato
      printf("%s", Messaggi[i].mittente.email);
}
    \end{clangsmall}
	
\end{Answer}
%\begin{Exercise}[label={ex:complemento}]
		\begin{itemize}
				\item[(a)] Si dica qual \`e l'intervallo di valori interi rappresentabile con la codifica in complemento a due a 9 bit.
				\item[(b)] Con riferimento a tale codifica indicare, giustificando brevemente le risposte, quali delle seguenti operazioni possono essere effettuate correttamente:
				\begin{itemize}
						\item $-254 - 255$
						\item $+254 - 253$
						\item $-18 + 236$
						\item $+217 + 182$
				\end{itemize}
				\item[(c)] Mostrare in dettaglio come avviene il calcolo delle operazioni (i) e (ii), evidenziando il bit di riporto e il bit di overflow cos\`i ottenuti.
		\end{itemize}

\end{Exercise}

\begin{Answer}[ref={ex:complemento}]
		\begin{itemize}
				\item[(a)] I valori rappresentabili vanno da $- 2^{m-1} = -256$ a $2^{m-1}-1 = +255$ compresi.
				\item[(b)] Le soluzioni delle operazioni sono:
				\begin{itemize}
						\item $-254 - 255$ NO si ottiene un valore negativo troppo grande in valore assoluto
						\item $+254 - 253$ SI si ottiene un valore piccolo in valore assoluto
						\item $- 18 + 236$ SI si ottiene un valore positivo, grande in valore assoluto ma nei limiti
						\item $+217 + 182$ NO si ottiene un valore positivo troppo grande in valore assoluto
						\end{itemize}
				\item[(c)]
						\item In complemento a 2 a $m = 9$ bit, $(-254)_{10} = (100000010)_{CP2}$ (perch\`e $(254)_{10} = (011111110)_2$). Da questo possiamo ricavare che $(-255)_{10} = (100000001)_{CP2}$ essendo $-255 = -254 - 1$. La somma diventa:\\
						\begin{table}[H]
						\centering
						\begin{tabular}{ccccccccccc|c}
								&& (1)\\ 
								&& 1 & 0 & 0 & 0 & 0 & 0 & 0 & 1 &0 &+\\
								&& 1 & 0 & 0 & 0 & 0 & 0 & 0 & 0 &1 &=\\
								\hline
								[1]&(1)& 0 & 0 & 0 & 0 & 0 & 0 & 0 & 1 &1 &\\
						\end{tabular}
						\end{table}
						\item $254$ in complemento a due \`e $(011111110)_{CP2} = (011111110)_2$ essendo positivo. Invece $-253 = -254 +1 = (100000011)_{CP2}$, quindi:\\
						\begin{table}[H]
						\centering
						\begin{tabular}{ccccccccccc|c}
								&      &(1)&(1)&(1)&(1)&(1)&(1)&(1)&(1)&  & \\
								&      & 0 & 1 & 1 & 1 & 1 & 1 & 1 & 1 &0 &+\\
								&      & 1 & 0 & 0 & 0 & 0 & 0 & 0 & 1 &1 &=\\
								\hline
								[0]&(1)& 0 & 0 & 0 & 0 & 0 & 0 & 0 & 0 &1 &\\
						\end{tabular}
						\end{table}
						ignorando il carry abbiamo il risultato esatto $(000000001)_{CP2} = (1)_{10}$.
				\end{itemize}
\end{Answer}
%\begin{Exercise}[label={ex:cp2}]
	\begin{enumerate}
		\item Si dica quale dei seguenti numeri \`e il maggiore, motivando con calcoli e ragionamenti la risposta:
			\begin{itemize}
				\item[A.] 34 (numero in base 10);
				\item[B.] 110110 (numero naturale in binario);
				\item[C.] 57 (numero in base 10);
				\item[D.] (numero in Complemento a 2);
				\item[E.] (numero in Complemento a 2).
			\end{itemize}
		\item Si indichi quindi il numero \emph{n} di bit che permette di codificare tutti i numeri A - E in Complemento a 2 e li si codifichi tutti in Complemento a 2 usando \emph{n} bit, riportando i passaggi fondamentali dei calcoli.
		\item Si eseguano (riportando i calcoli) le seguenti operazioni utilizzando \emph{n} bit e si indichino eventuali bit di carry (riporto) e overflow. Si segnalino i casi in cui l'operazione non pu\'o essere eseguita utilizzando gli \emph{n} bit identificati nel punto 2:
			\begin{itemize}
				\item[I.] A - B;
				\item[II.] (C + D) + E;
				\item[III.] 2A + E.
			\end{itemize}
	\end{enumerate}

\end{Exercise}

\begin{Answer}[ref={ex:cp2}]
	\begin{enumerate}
		\item Tra i numeri rappresentati in base 10 possiamo immediatamente decidere che il numero maggiore \`e 57. Il numero 1111111 \`e 		chiaramente negativo avendo il bit pi\'u significativo pari a 1, quindi \`e certamente minore di 57. Possiamo convertire in base 10 anche gli altri due numeri per prendere la decisione finale:
		\begin{itemize}
			\item 110110:  32+16+4+2 = 54;
			\item 0011011:  16+8+2+1 = 27;
		\end{itemize}
si noti che 0011011 \`e in complemento a 2 (CP2) ma \`e un numero positivo (il bit pi\'u significativo \`e pari a 0), di conseguenza, \`e possibile applicare la regola di conversione che vale anche per i numeri naturali.
Possiamo quindi concludere che il numero maggiore \`e 57.
	\item Il numero maggiore, 57, necessita di 7 bit per la conversione, infatti n deve essere tale che $57 \leq 2^{n - 1}$. Per n = 7 si ha che $2^{n - 1}$ \`e pari a 64. Se avessimo scelto n pari a 6, avremmo avuto che $2^{n - 1} = 32$, che è chiaramente minore di 57. Anche l'unico numero negativo \`e espresso su 7 cifre binarie, quindi 7 \`e il numero di bit adeguato a convertire tutti i numeri considerati. Convertendo tutti i numeri in CP2 a 7 bit si ha:
		\begin{itemize}
			\item[A.] 34 in base 10 $\Rightarrow$ 0100010 in CP2 su 7 bit;
			\item[B.] 110110 in binario $\Rightarrow$ 0110110 in CP2 su 7 bit;
			\item[C.] 57 in base 10 $\Rightarrow$ 0111001 in CP2 su 7 bit;
			\item[D.] 1111111 in CP2 su 7 bit;
			\item[E.] 0011011 in CP2 su 7 bit.
		\end{itemize}
	\item 
	\begin{itemize}
		\item[I.] Prima calcolo $- B$ = 10001010 ed ho:
		
			0 1 0 0 0 1 0 + 1 0 0 1 01 0 = [ 0 ] 1 1 0 1 1 0 0;

		\item[II.] Eseguo le due somme:

			0 1 1 1 0 0 1 + 1 1 1 1 1 1 1 = 0 1 1 1 0 0 0 + 0 0 1 1 0 1 1 = [1] 1 0 1 0 0 1 1;

		In questo caso abbiamo overflow.
		\item[III.] Questa operazione non \`e fattibile in quanto 2A = 68 che darebbe gi\'a overflow.
	\end{itemize}
	\end{enumerate}

\end{Answer}
%\input{ex/ex_ieee_standard}
%\input{ex/ex_binary}
%\input{ex/ex_ieee_standard_c}

\clearpage
\section*{Soluzioni}
\shipoutAnswer

\end{document}